Text-based machine-learning and Transformer models were evaluated to examine the contribution of linguistic information to Urdu SER. Table~\ref{tab:text_based_models} summarizes their performance. Among the conventional machine-learning approaches, TF-IDF with Linear SVM achieved the highest performance, obtaining 84.00\% accuracy and 83.95\% macro-F1. TF-IDF with Multinomial Naive Bayes also demonstrated competitive performance, achieving 77.14\% accuracy and 77.05\% macro-F1, while TF-IDF with Logistic Regression reached 75.71\% accuracy and 75.60\% macro-F1.

Among the Transformer-based baselines, Distil-mBERT achieved the highest accuracy of 78.57\% with a macro-F1 of 78.31\%, followed by XLM-RoBERTa-large with 77.43\% accuracy and 77.44\% macro-F1. The remaining multilingual Transformer models, including mBERT and mDeBERTa-v3-base, provided competitive but lower performance in the seven-class Urdu SER setting.

The strong performance of TF-IDF with Linear SVM indicates that lexical patterns and discriminative word-level information remain highly informative for emotion classification in the Urdu transcripts. Although Transformer models provide contextualized representations, their performance did not surpass the optimized classical TF-IDF-based approach on this dataset, suggesting that dataset characteristics and linguistic patterns play an important role in Urdu emotion recognition.

The proposed multimodal fusion framework achieved 89.14\% accuracy and 89.10\% macro-F1, substantially outperforming all text-only baselines. This improvement demonstrates that while linguistic information captures a significant portion of emotion-related cues, the integration of complementary acoustic characteristics provides additional emotional information that is not fully represented in textual transcripts alone.
